% Standalone problem document, generated from the adjacent README.md.
% Compile with XeLaTeX. Mathematical content is maintained in Markdown.
\documentclass[11pt,a4paper]{article}
\usepackage[fontsize=10.5pt]{scrextend}
\usepackage[margin=22mm,headheight=15pt,headsep=9mm,footskip=11mm]{geometry}
\usepackage{fontspec}
\setmainfont{texgyrepagella}[Extension=.otf,UprightFont=*-regular,BoldFont=*-bold,ItalicFont=*-italic,BoldItalicFont=*-bolditalic]
\setsansfont{texgyreheros}[Extension=.otf,UprightFont=*-regular,BoldFont=*-bold,ItalicFont=*-italic,BoldItalicFont=*-bolditalic]
\setmonofont{texgyrecursor}[Extension=.otf,UprightFont=*-regular,BoldFont=*-bold,ItalicFont=*-italic,BoldItalicFont=*-bolditalic,Scale=MatchLowercase]
\usepackage{amsmath,amssymb}
\usepackage{unicode-math}
\setmathfont{texgyrepagella-math.otf}
\usepackage{xcolor,microtype,enumitem,needspace,fancyhdr}
\usepackage{bookmark}
\definecolor{ink}{HTML}{17344B}
\definecolor{accent}{HTML}{197D80}
\definecolor{muted}{HTML}{596773}
\hypersetup{colorlinks=true,linkcolor=accent,urlcolor=accent,pdftitle={RE-06 - Nonadaptive
queries for finite-family matrix
approximation},pdfauthor={Open Problems in Numerical Linear Algebra}}
\urlstyle{same}
\setlength{\parindent}{0pt}
\setlength{\parskip}{5pt plus 1pt minus 1pt}
\linespread{1.04}
\setlength{\emergencystretch}{3em}
\widowpenalty=10000
\clubpenalty=10000
\displaywidowpenalty=10000
\allowdisplaybreaks[1]
\setlist{nosep,leftmargin=1.5em}
\providecommand{\tightlist}{\setlength{\itemsep}{0pt}\setlength{\parskip}{0pt}}
\providecommand{\pandocbounded}[1]{#1}
\pagestyle{fancy}
\fancyhf{}
\fancyhead[L]{\sffamily\footnotesize\color{muted}OPEN PROBLEMS IN NUMERICAL LINEAR ALGEBRA}
\fancyhead[R]{\sffamily\bfseries\footnotesize\color{accent}RE-06}
\fancyfoot[L]{\parbox[b]{0.9\textwidth}{\sffamily\fontsize{8}{10}\selectfont\color{muted}Editorial ratings; a bounded search cannot certify that a problem remains unsolved.\\Literature check: 2026-09-10}}
\fancyfoot[R]{\sffamily\footnotesize\color{muted}\thepage}
\renewcommand{\headrulewidth}{0.3pt}
\renewcommand{\headrule}{\hbox to\headwidth{\color{accent}\leaders\hrule height \headrulewidth\hfill}}
\makeatletter
\renewcommand{\section}{\Needspace{5\baselineskip}\@startsection{section}{1}{0pt}{12pt plus 2pt minus 2pt}{4pt}{\sffamily\bfseries\large\color{ink}}}
\renewcommand{\subsection}{\Needspace{4\baselineskip}\@startsection{subsection}{2}{0pt}{9pt plus 1pt minus 1pt}{3pt}{\sffamily\bfseries\normalsize\color{ink}}}
\makeatother
\setcounter{secnumdepth}{0}
\begin{document}
{\sffamily\small\color{muted}Randomized and low-rank approximation\par}
\vspace{3pt}
{\raggedright\hyphenpenalty=10000\exhyphenpenalty=10000\sffamily\bfseries\fontsize{21}{24}\selectfont\color{ink}Nonadaptive
queries for finite-family matrix approximation\par}
\vspace{7pt}
{\sffamily\small\textbf{Difficulty:} challenging\quad\textbf{Importance:} interesting
to the community\par}
{\sffamily\small\bfseries\color{accent}Status: Open\par}
\vspace{4pt}
{\color{accent}\hrule height 0.8pt}
\vspace{6pt}
\textbf{Rating rationale:} Challenging because all query choices must
precede information about the target; community impact is parallel and
pass-efficient structured matrix learning.

\textbf{Area:} nonadaptive matrix sketching

\subsection{Context and notation}\label{context-and-notation}

Use exact real arithmetic, with comparisons, standard Gaussian sampling,
and exact SVDs available as primitives; charge an SVD of an
\(a\times b\) matrix \(O(ab\min(a,b))\) operations. This states the
idealized arithmetic model used here, rather than a finite precision or
bit complexity claim. A matrix--vector query returns either \(Av\) or
\(A^\mathsf Tv\) for one chosen real vector \(v\); both types count
toward the total. Randomized guarantees are for every fixed input, with
probability or expectation over the algorithm's randomness.

\subsection{Problem statement}\label{problem-statement}

Let \(n\ge1\), let \(\mathcal F\subset\mathbb R^{n\times n}\) be
explicitly given with \(M=|\mathcal F|\ge2\), and let
\(A\in\mathbb R^{n\times n}\) be accessible only through matrix--vector
queries. Set \(t=\log(2M)\).

\subsubsection{Question}\label{question}

Do absolute constants \(C>0\), integer \(b\ge0\), and a uniform
randomized algorithm exist that, for every \(0<\varepsilon<1/2\), choose
all query vectors and all choices between \(A\) and \(A^\mathsf T\)
before receiving any oracle answers, make at most

\[
C\sqrt t\,\varepsilon^{-2}
\bigl[1+\log(2+t)+\log(1/\varepsilon)\bigr]^b
\]

queries, and return \(B\in\mathcal F\) such that

\[
\Pr\!\left[\|A-B\|_F\le(3+\varepsilon)
\min_{D\in\mathcal F}\|A-D\|_F\right]\ge0.99?
\]

Query choices may depend on \(\mathcal F,\varepsilon\) and randomness.
Only oracle calls are charged; candidate processing and processing of
the answers are unrestricted.

\newpage
\subsection{References}\label{references}

Amsel et al., \href{https://arxiv.org/html/2507.19290v2}{\emph{Query
Efficient Structured Matrix Learning}}, §5, adaptivity question.
Christopher Musco,
\href{https://app.icerm.brown.edu/assets/568/10574/10574_5858_Musco_020420261630_Slides.pdf}{\emph{Structured
Matrix Learning from Matrix-Vector Products}}, ICERM, February 4, 2026,
numbered slide 23.

\subsection{Status evidence}\label{status-evidence}

Both sources explicitly ask whether adaptivity is necessary; the paper's
August 21, 2026 revision retains the question. Searches for ``finite
matrix approximation adaptivity 2026'' and the paper identifier with
``non-adaptive'' found no resolution. The abstract's linked
finite-family improvement, Theorem 1.2, uses two rounds after receiving
a constant-factor warm start: its left query vectors depend on the first
round's answers. It does not supply the required nonadaptive algorithm.

\subsection{Audit --- 2026-09-10}\label{audit-2026-09-10}

Rechecked \href{https://arxiv.org/html/2507.19290v2}{Amsel et al., §5}
and the \href{https://proceedings.mlr.press/v336/amsel26a.html}{COLT
2026 record}. The adaptivity question remains. Nonadaptive
structured-learning searches found no resolution; the linked
finite-family relative-error improvement still uses responses to choose
later queries.
\end{document}
